\section{Introduction}

Large language models (LLMs) have become a fundamental component of modern AI systems, with most implementations heavily optimized for GPU-based platforms such as NVIDIA A100. However, the increasing demand for efficient inference and cost-effective deployment has led to the emergence of alternative AI accelerators, including the Tenstorrent Wormhole architecture.

This project explores the problem of porting Granite-4.0-h small and tiny models — IBM's hybrid language models combining recurrent and attention layers with sparse expert routing — from the standard HuggingFace/PyTorch stack to the Tenstorrent Wormhole Galaxy 32 multi-chip system. On a GPU, hundreds of compute operations within a single forward pass are managed efficiently by mature GPU runtimes. On Wormhole, each such operation must be explicitly sent from the host CPU to the device one at a time, and each transfer takes a fixed amount of time regardless of how much computation that operation actually does. During token generation — where the model processes just one token at a time — this per-operation cost accumulates and becomes the dominant bottleneck. Section~\ref{sec:arch} introduces the hardware in detail.

The key research questions addressed in this work are:
\begin{itemize}
\item What numerical and architectural challenges arise when porting a hybrid recurrent-attention model with sparse expert routing from HuggingFace to Wormhole?
\item How does the per-operation communication cost between the host CPU and the device limit token generation speed, and what optimisations address it?
\item When model weights are split across multiple chips and results must be merged after each layer, what additional constraints does this introduce for efficient token generation?
\end{itemize}

The contributions of this project are:
\begin{itemize}
\item Complete working implementations of Granite-4.0-h-tiny (4 Wormhole chips) and Granite-4.0-h-small (8 chips), with token generation throughput benchmarked against NVIDIA A100 and CPU baselines.
\item A systematic characterisation of the numerical, architectural, and toolchain challenges encountered during the port (Section~\ref{sec:challenges}).
\item A custom on-chip kernel for the recurrent state-update step that reduces memory traffic and the number of host-to-device commands per token.
\item A working implementation of whole-model operation recording and replay — a technique where all 500+ operations per token are recorded once into a device-side buffer and replayed with a single command on every subsequent token. We also identify and fix a subtle bug where recording inside the generation loop causes every output token to be shifted by one position.
\item Empirical token generation throughput measurements and per-component timing breakdowns.
\end{itemize}
